%! Displaying and Analyzing Categorical Data — Unit 3, Lesson 3.1 — Slides (Beamer)
% GRADUAL RELEASE deck, 5/45/5/5: title → targets → warm-up → hook (dark, UNRESOLVED) → I DO
% divider (dark) → one frame per notes row, crux frames flagged \sectionlabel[redacc]{} → Guided
% Practice (a LIVE surface, un-answered) → spoken debrief with the cold check → You do: THE
% HOMEWORK, assigned (not started in class) → close (dark). Standards: PS.DS.4a-d.
\documentclass[aspectratio=169, 11pt]{beamer}
\usepackage{saar-beamer}   % provides \ceruleanheader{} and \sectionlabel[color]{}

\newcommand{\redink}[1]{{\color{redacc}\bfseries #1}}

\begin{document}

% TITLE SLIDE (CourseName is not defined in beamer, so the name is literal)
{
\setbeamercolor{background canvas}{bg=cerulean}
\begin{frame}[plain]
  \vspace{2.8cm}\hspace{0.55cm}%
  \begin{minipage}{0.9\textwidth}
    {\color{goldacc}\bfseries\small LESSON 3.1}\\[0.5cm]
    {\color{white}\bfseries\LARGE Displaying and Analyzing Categorical Data}\\[1.2cm]
    {\color{frostmid}\small Statistical Analysis \& Algebraic Reasoning~$\cdot$~Unit 3~$\cdot$~PS.DS.4a--d}
  \end{minipage}
\end{frame}
}

% ── LEARNING TARGETS ─────────────────────────────────────────────────────────
\begin{frame}[t, plain]
  \ceruleanheader{Today's targets}
  \sectionlabel{Lesson 3.1}
  By the end of today, I can\ldots
  \begin{itemize}\setlength{\itemsep}{3pt}
    \item read a bar graph as a frequency or a \textbf{relative frequency}.
    \item compare groups with \textbf{side-by-side} and \textbf{segmented} bar graphs.
    \item build a \textbf{two-way table} and read its \textbf{marginal frequencies}.
    \item find \textbf{conditional relative frequencies} and name the whole.
    \item compare groups of different sizes with percents, not raw counts.
  \end{itemize}
  \begin{block}{How today works}
    Warm-up $\to$ \textbf{notes together} $\to$ \textbf{one survey we work as a class} $\to$
    debrief $\to$ \textbf{the homework, assigned today and done outside class}.
  \end{block}
\end{frame}

% ── WARM-UP ──────────────────────────────────────────────────────────────────
\begin{frame}[t, plain]
  \ceruleanheader{Warm-Up --- 5 minutes, silent}
  \sectionlabel{Unit 3 starts here $\cdot$ groups and chances}
  \vspace{0.2cm}
  \begin{center}
    \begin{minipage}{0.86\textwidth}
      \begin{enumerate}\setlength{\itemsep}{7pt}
        \item Find 47\% of 200, 45\% of 120, and 50\% of 80.
        \item Write as percents: $\dfrac{94}{200}$, $\dfrac{54}{120}$, $\dfrac{40}{80}$. Which of
              the last two has the larger top number? Which is the larger percent?
        \item Bus 94, Car 80, Walk 26, Total 200. How many come by car? What percent walk?
      \end{enumerate}
    \end{minipage}
  \end{center}
  \vspace{0.2cm}
  \begin{center}
    {\color{cerulean}\bfseries Keep the page --- every number comes back in the notes.}
  \end{center}
\end{frame}

% ── HOOK — leave it UNRESOLVED; problem 9 settles it ─────────────────────────
{
\setbeamercolor{background canvas}{bg=cerulean}
\begin{frame}[plain]
  \vspace{1.2cm}
  \begin{center}
    {\color{frostmid}\bfseries\small THE OAKDALE HIGH NEWSLETTER}\\[0.7cm]
    {\color{white}\bfseries\Large ``54 ninth graders ride the bus and only 40 twelfth graders do,
    so ninth graders are more likely to ride the bus.''}\\[0.8cm]
    {\color{goldacc}\bfseries\large Agree, disagree, or can't tell yet?}\\[0.5cm]
    {\color{frostmid}\small Vote, and write one line of why. We vote again at problem 9.}
  \end{center}
\end{frame}
}

% ── I DO ─────────────────────────────────────────────────────────────────────
{
\setbeamercolor{background canvas}{bg=cerulean}
\begin{frame}[plain]
  \vspace{1.8cm}\hspace{0.8cm}%
  \begin{minipage}{0.86\textwidth}
    {\color{goldacc}\bfseries\small GUIDED NOTES \ \textbullet\ \ I DO}\\[0.7cm]
    {\color{white}\bfseries\Large Open to your Guided Notes page. We work the table row by
    row --- fill each term as we name it, not before.}\\[0.9cm]
    {\color{frostmid}\small Four terms go in the vocabulary box today. Write each one the
    moment we use it.}
  \end{minipage}
\end{frame}
}

% ── ROW 1. BAR GRAPHS ────────────────────────────────────────────────────────
\begin{frame}[t, plain]
  \ceruleanheader{One variable, two axes}
  \sectionlabel{notes row 1 $\cdot$ Bar Graphs $\cdot$ problems 1--2}
  \begin{center}
  \begin{tikzpicture}[x=1cm, y=0.03cm]
    \foreach \v in {20,40,...,100} { \draw[linegray] (0,\v) -- (4.6,\v); }
    \foreach \v in {0,20,...,100} { \node[left, font=\scriptsize] at (0,\v) {\v}; }
    \draw[thick] (0,0) -- (0,104);  \draw[thick] (0,0) -- (4.6,0);
    \foreach \x/\h/\f/\n in {0.45/94/cerulean/Bus, 1.9/80/frostmid/Car, 3.35/26/goldbg/Walk} {
      \fill[\f, draw=charcoal] (\x,0) rectangle ({\x+0.8},\h);
      \node[above, font=\scriptsize\bfseries] at ({\x+0.4},\h) {\h};
      \node[below, font=\scriptsize] at ({\x+0.4},0) {\n};
    }
    \node[font=\footnotesize\bfseries] at (2.3,118) {Frequency};
    \begin{scope}[xshift=6.7cm]
    \foreach \v in {10,20,...,50} { \draw[linegray] (0,{\v*2}) -- (4.6,{\v*2}); }
    \foreach \v in {0,10,...,50} { \node[left, font=\scriptsize] at (0,{\v*2}) {\v\%}; }
    \draw[thick] (0,0) -- (0,104);  \draw[thick] (0,0) -- (4.6,0);
    \foreach \x/\h/\f/\n/\p in {0.45/94/cerulean/Bus/47, 1.9/80/frostmid/Car/40, 3.35/26/goldbg/Walk/13} {
      \fill[\f, draw=charcoal] (\x,0) rectangle ({\x+0.8},\h);
      \node[above, font=\scriptsize\bfseries] at ({\x+0.4},\h) {\p\%};
      \node[below, font=\scriptsize] at ({\x+0.4},0) {\n};
    }
    \node[font=\footnotesize\bfseries] at (2.3,118) {Relative frequency};
    \end{scope}
  \end{tikzpicture}
  \end{center}
  \begin{itemize}\setlength{\itemsep}{3pt}
    \item Problem 1: \textbf{94} of the 200 surveyed students ride the bus.
    \item Problem 2: Car $= 80/200 = \mathbf{40\%}$. \redink{Relative frequency = count $\div$ total.}
  \end{itemize}
\end{frame}

% ── ROW 2. TWO KINDS OF BARS ─────────────────────────────────────────────────
\begin{frame}[t, plain]
  \ceruleanheader{Two ways to compare the grades}
  \sectionlabel{notes row 2 $\cdot$ Two Kinds of Bars $\cdot$ problems 3--4 $\cdot$ both graphs are counts}
  \begin{center}
  \begin{tikzpicture}[x=1cm, y=0.05cm]
    \foreach \v in {20,40,60} { \draw[linegray] (0,\v) -- (5.2,\v); }
    \foreach \v in {0,20,40,60} { \node[left, font=\scriptsize] at (0,\v) {\v}; }
    \draw[thick] (0,0) -- (0,63);  \draw[thick] (0,0) -- (5.2,0);
    \foreach \c/\a/\b/\n in {0.3/54/40/Bus, 2.0/48/32/Car, 3.7/18/8/Walk} {
      \fill[cerulean, draw=charcoal] (\c,0) rectangle ({\c+0.6},\a);
      \fill[goldbg, draw=charcoal] ({\c+0.6},0) rectangle ({\c+1.2},\b);
      \node[above, font=\scriptsize\bfseries] at ({\c+0.3},\a) {\a};
      \node[above, font=\scriptsize\bfseries] at ({\c+0.9},\b) {\b};
      \node[below, font=\scriptsize] at ({\c+0.6},0) {\n};
    }
    \node[font=\scriptsize\bfseries] at (2.6,72) {Side-by-side (dark: Grade 9)};
    \begin{scope}[xshift=7cm, y=0.0275cm]
    \foreach \v in {40,80,120} { \draw[linegray] (0,\v) -- (4.8,\v); }
    \foreach \v in {0,40,80,120} { \node[left, font=\scriptsize] at (0,\v) {\v}; }
    \draw[thick] (0,0) -- (0,122);  \draw[thick] (0,0) -- (4.8,0);
    \foreach \x/\a/\b/\t/\g in {0.5/54/102/120/Grade 9, 2.9/40/72/80/Grade 12} {
      \fill[cerulean, draw=charcoal] (\x,0) rectangle ({\x+1},\a);
      \fill[frostmid, draw=charcoal] (\x,\a) rectangle ({\x+1},\b);
      \fill[goldbg, draw=charcoal] (\x,\b) rectangle ({\x+1},\t);
      \node[right, font=\scriptsize\bfseries, inner sep=1.5pt] at ({\x+1},\a) {\a};
      \node[right, font=\scriptsize\bfseries, inner sep=1.5pt] at ({\x+1},\b) {\b};
      \node[above, font=\scriptsize\bfseries, inner sep=1.5pt] at ({\x+0.5},\t) {\t};
      \node[below, font=\scriptsize] at ({\x+0.5},0) {\g};
    }
    \node[font=\scriptsize\bfseries] at (2.4,143) {Segmented (running totals)};
    \end{scope}
  \end{tikzpicture}
  \end{center}
  \begin{itemize}\setlength{\itemsep}{3pt}
    \item Problem 3: the largest gap is Car, $48 - 32 = \mathbf{16}$ students.
    \item Problem 4: twelfth graders by car: $72 - 40 = \mathbf{32}$. \redink{72 is a running total.}
  \end{itemize}
\end{frame}

% ── ROW 3. TWO-WAY TABLE ─────────────────────────────────────────────────────
\begin{frame}[t, plain]
  \ceruleanheader{From a summary to a two-way table}
  \sectionlabel{notes row 3 $\cdot$ Two-Way Tables $\cdot$ problems 5--7}
  \begin{columns}[t]
    \begin{column}{0.54\textwidth}
      \renewcommand{\arraystretch}{1.4}
      \begin{tabular}[t]{l c c c c}
      \hline
       & \textbf{Bus} & \textbf{Car} & \textbf{Walk} & \textbf{Total} \\
      \hline
      Grade 9 & 54 & 48 & \redink{18} & 120 \\
      Grade 12 & 40 & 32 & \redink{8} & 80 \\
      \textbf{Total} & \textbf{94} & \textbf{80} & \textbf{26} & \textbf{200} \\
      \hline
      \end{tabular}
    \end{column}
    \begin{column}{0.42\textwidth}
      \begin{itemize}\setlength{\itemsep}{5pt}
        \item The summary leaves out Walk: \textbf{subtract}.
        \item Inside cell: a \textbf{joint} count. Problem 6: \textbf{8} twelfth graders walk.
        \item A Total: a \textbf{marginal frequency}. Problem 7: \textbf{80} come by car.
      \end{itemize}
    \end{column}
  \end{columns}
\end{frame}

% ── ROW 4. WHICH WHOLE ───────────────────────────────────────────────────────
\begin{frame}[t, plain]
  \ceruleanheader{Which whole?}
  \sectionlabel{notes row 4 $\cdot$ Which Whole? $\cdot$ problem 8, then vote again at problem 9}
  \begin{itemize}\setlength{\itemsep}{6pt}
    \item \textbf{Marginal relative frequency:} a Total $\div$ the whole table.
    \item \textbf{Conditional relative frequency:} an inside cell $\div$ the total of
          \emph{one} group.
    \item Problem 8: $94/200 = \mathbf{47\%}$ of the surveyed students ride the bus.
  \end{itemize}
  \begin{block}{The test}
    In ``Of the [group], what percent \ldots?'' \redink{the words after \emph{of} name the whole}
    --- and that total is the denominator.
  \end{block}
  \vspace{6pt}
  \begin{center}
    {\color{cerulean}\bfseries Problem 9: vote again before anyone divides.}
  \end{center}
\end{frame}

% ── THE CRUX ─────────────────────────────────────────────────────────────────
\begin{frame}[t, plain]
  \ceruleanheader{More riders. Smaller percent.}
  \sectionlabel[redacc]{notes row 4 $\cdot$ problem 9 $\cdot$ the whole point of today}
  \begin{center}
  \begin{tikzpicture}[x=0.1cm, y=1cm]
    \foreach \v in {0,20,...,100} {
      \draw[linegray] (\v,0) -- (\v,1.75);
      \node[below, font=\scriptsize] at (\v,0) {\v\%}; }
    \draw[thick] (0,0) -- (100,0);
    \foreach \y/\g/\a/\b/\c in {1.05/Grade 9/45/40/15, 0.2/Grade 12/50/40/10} {
      \fill[cerulean, draw=charcoal] (0,\y) rectangle (\a,\y+0.55);
      \fill[frostmid, draw=charcoal] (\a,\y) rectangle ({\a+\b},\y+0.55);
      \fill[goldbg, draw=charcoal] ({\a+\b},\y) rectangle (100,\y+0.55);
      \node[font=\footnotesize\bfseries, text=white] at ({\a/2},\y+0.275) {Bus \a\%};
      \node[font=\footnotesize\bfseries] at ({\a+\b/2},\y+0.275) {Car \b\%};
      \node[font=\scriptsize\bfseries] at ({\a+\b+\c/2},\y+0.275) {\c\%};
      \node[anchor=east, font=\scriptsize\bfseries] at (-1,\y+0.275) {\g};
    }
  \end{tikzpicture}
  \end{center}
  \begin{itemize}\setlength{\itemsep}{2pt}
    \item Ninth graders: $54/120 =$ \redink{45\%}. Twelfth graders: $40/80 =$ \redink{50\%}.
  \end{itemize}
  \begin{block}{Counts say how many. Conditional percents say how likely.}
    To compare groups of different sizes, compare conditional percents. The newsletter is wrong:
    \redink{twelfth graders are more likely to ride.}
  \end{block}
\end{frame}

% ── THE TRAP ─────────────────────────────────────────────────────────────────
\begin{frame}[t, plain]
  \ceruleanheader{The same 54, two different wholes}
  \sectionlabel[redacc]{notes row 4 $\cdot$ problems 10--11}
  \vspace{4pt}
  \begin{columns}[t]
    \begin{column}{0.47\textwidth}
      \begin{block}{Of ninth graders, what percent ride?}
        $54/120 = 45\%$. The whole is the \textbf{120 ninth graders}.
      \end{block}
    \end{column}
    \begin{column}{0.47\textwidth}
      \begin{block}{Of bus riders, what percent are ninth graders?}
        $54/94 \approx 57\%$. The whole is the \textbf{94 bus riders}.
      \end{block}
    \end{column}
  \end{columns}
  \vspace{8pt}
  \begin{itemize}\setlength{\itemsep}{4pt}
    \item Problem 10: Dani's 57\% is a percent of \redink{bus riders}, not of ninth graders.
    \item Problem 11: Car is $48/120 = 40\%$ and $32/80 = 40\%$. \textbf{Different counts, the
          same likelihood.}
  \end{itemize}
\end{frame}

% ── WE DO — a LIVE surface: task and data only, no answers ───────────────────
\begin{frame}[t, plain]
  \ceruleanheader{Guided Practice: Streaming Choices}
  \sectionlabel[goldacc]{We do $\cdot$ problems 12--15 $\cdot$ you hold the pen}
  \begin{columns}[t]
    \begin{column}{0.5\textwidth}
      {\small\renewcommand{\arraystretch}{1.3}%
      \begin{tabular}{l c c c c}
      \hline
       & \textbf{Star} & \textbf{Reel} & \textbf{Pulse} & \textbf{Total} \\
      \hline
      Under 30 & 45 & 30 & 25 & 100 \\
      30 and over & 60 & 66 & 24 & 150 \\
      \textbf{Total} & 105 & 96 & 49 & 250 \\
      \hline
      \end{tabular}}\\[6pt]
      {\footnotesize Fairview Community Library, 250 patrons.}\\[8pt]
      {\color{cerulean}\bfseries Before you divide: of whom?}
    \end{column}
    \begin{column}{0.46\textwidth}
      {\small
      \begin{itemize}\setlength{\itemsep}{3pt}
        \item[12.] Percent of all patrons who prefer Pulse? What kind?
        \item[13.] Of patrons under 30, percent who prefer Pulse? What kind?
        \item[14.] Ty: ``$25/49 \approx 51\%$, so 51\% of patrons under 30 prefer Pulse.'' What
              whole did he use?
        \item[15.] Flyer: ``More patrons 30 and over chose StarView (60 to 45), so they are more
              likely to prefer it.'' Check it.
      \end{itemize}}
    \end{column}
  \end{columns}
\end{frame}

% ── DEBRIEF ──────────────────────────────────────────────────────────────────
\begin{frame}[t, plain]
  \ceruleanheader{Debrief: of whom?}
  \sectionlabel{5 minutes $\cdot$ pens down, papers up --- we say these aloud}
  \vspace{4pt}
  \begin{enumerate}\setlength{\itemsep}{5pt}
    \item Problem 15: 40\% of patrons 30 and over, 45\% of patrons under 30. Say the sentence ---
          \emph{counts say how many; conditional percents say how likely.}
  \end{enumerate}
  \vspace{6pt}
  \begin{block}{Cold check --- hands down, thirty seconds}
    Thirty of \textbf{60} juniors and \textbf{40} of \textbf{100} seniors signed up for a field
    trip. Morgan says seniors are more likely to sign up. \textbf{Is she right? Give both
    percents.}
  \end{block}
\end{frame}

% ── YOU DO — the homework is the individual practice, assigned today ─────────
\begin{frame}[t, plain]
  \ceruleanheader{You do: the homework}
  \sectionlabel[goldacc]{Assigned today $\cdot$ on your own, outside class $\cdot$ due after two study halls}
  \begin{itemize}\setlength{\itemsep}{5pt}
    \item \textbf{Pet adoptions at an animal shelter} --- dog or cat, by adopter age.
    \item Finish a two-way table from a summary; interpret a percent.
    \item Two questions about the same count, over two different wholes.
    \item A claim built on counts, and a choice between two displays.
  \end{itemize}
  \begin{block}{DeltaMath (two-way tables and their percents) + turn in packet items 2, 5, 6. Scored 1 to 5, due at the start of the first class after two study halls.}
    Do item 5 first if you are short on time --- it shows whether today landed. Your notes' rows
    3 and 4 and the \emph{Keep in Mind} box are your reference.
  \end{block}
\end{frame}

% ── CLOSE ────────────────────────────────────────────────────────────────────
{
\setbeamercolor{background canvas}{bg=cerulean}
\begin{frame}[plain]
  \vspace{1.5cm}\hspace{0.8cm}%
  \begin{minipage}{0.86\textwidth}
    {\color{goldacc}\bfseries\small BEFORE YOU GO}\\[0.7cm]
    {\color{white}\bfseries\Large Counts say how many. Conditional percents say how likely.
    Before you divide, ask: of whom?}\\[0.9cm]
    {\color{frostmid}\small Next: Lesson 3.2, Probability Models --- today's relative frequency
    is the first step toward a probability.}
  \end{minipage}
\end{frame}
}

\end{document}
